\section{Introduction}
Autonomous agents utilizing Large Language Models (LLMs) offer promising opportunities to enhance and replicate human workflows. 
%
In real-world applications, however, existing systems~\citep{park2023generative, zhuge2023mindstorms,cai2023large,wang2023unleashing,li2023camel,du2023improving,liang2023encouraging,hao2023chatllm, zhou2023agents} tend to oversimplify the complexities. 
%
They struggle to achieve effective, coherent, and accurate problem-solving processes, particularly when there is a need for meaningful collaborative interaction~\citep{chen2024s,zhang2023building,dong2023self,zhou2023webarena, qian2023communicative,tang2023medagents,hong2024data}.

Through extensive collaborative practice, humans have developed widely accepted Standardized Operating Procedures (SOPs) across various domains~\citep{belbin2012team,manifesto2001manifesto,demarco2013peopleware}.
%
These SOPs play a critical role in supporting task decomposition and effective coordination. 
%
Furthermore, SOPs outline the responsibilities of each team member, while establishing standards for intermediate outputs. 
%
Well-defined SOPs improve the consistent and accurate execution of tasks that align with defined roles and quality standards~\citep{belbin2012team, manifesto2001manifesto,demarco2013peopleware,10.1145/280765.280867}. 
%
For instance, in a software company, Product Managers analyze competition and user needs to create Product Requirements Documents (PRDs) using a standardized structure, to guide the developmental process. 
%

\begin{figure}
    \centering
    \vspace{-1cm}
    \includegraphics[width=\textwidth]{imgs/1-metagpt_overall_update.png}
    %\vspace{-0.4cm}
    \caption{
    \textbf{The software development SOPs between MetaGPT and real-world human teams.}
    %
    In software engineering, SOPs promote collaboration among various roles.
    %
    MetaGPT showcases its ability to decompose complex tasks into specific actionable procedures assigned to various roles (e.g., Product Manager, Architect, Engineer, etc.). 
    }
    \label{fig:SOP}
\end{figure}

Inspired by such ideas, we design a promising \underline{GPT}-based \underline{Meta}-Programming framework called {MetaGPT} that significantly benefits from SOPs. 
%
Unlike other works~\citep{li2023camel,qian2023communicative}, MetaGPT requires agents to generate structured outputs, such as high-quality requirements documents, design artifacts, flowcharts, and interface specifications. 
%
The use of intermediate structured outputs significantly increases the success rate of target code generation.
%
Because it helps maintain consistency in communication, minimizing ambiguities and errors during collaboration.
%
More graphically, in a company simulated by MetaGPT, all employees follow a strict and streamlined workflow, and all their handovers must comply with certain established standards. This reduces the risk of hallucinations caused by idle chatter between LLMs, particularly in role-playing frameworks, like: ``\textit{Hi, hello and how are you?''} -- Alice (Product Manager); ``\textit{Great! Have you had lunch?''} -- Bob (Architect).


Benefiting from SOPs, MetaGPT offers a promising approach to meta-programming. In this context, we adopt meta-programming\footnote{\href{https://en.wikipedia.org/w/index.php?title=Metaprogramming}{https://en.wikipedia.org/w/index.php?title=Metaprogramming}} as "programming to program", in contrast to the broader fields of meta learning and "learning to learn"~\citep{schmidhuber1987evolutionary,Schmidhuber:93selfreficann,Hochreiter:01meta,Schmidhuber:05gmai,finn2017model}. 

This notion of meta-programming also encompasses earlier efforts like CodeBERT~\citep{feng2020codebert} and recent projects such as CodeLlama~\citep{roziere2023code} and WizardCoder~\citep{luo2023wizardcoder}. 
However, MetaGPT stands out as a unique solution that allows for efficient meta-programming through a well-organized group of specialized agents. Each agent has a specific role and expertise, following some established standards. This allows for automatic requirement analysis, system design, code generation, modification, execution, and debugging during runtime, highlighting how agent-based techniques can enhance meta-programming.


To validate the design of MetaGPT, we use publicly available HumanEval~\citep{chen2021evaluating} and MBPP~\citep{austin2021program} for evaluations. 
%
Notably, in code generation benchmarks, MetaGPT achieves a new state-of-the-art (SoTA) with 85.9\% and 87.7\% in Pass@1.
%
When compared to other popular frameworks for creating complex software projects, such as AutoGPT~\citep{autogpt_github}, LangChain~\citep{Chase_LangChain_2022}, AgentVerse~\citep{agentverse_github}, and ChatDev~\citep{qian2023communicative}. MetaGPT also stands out in handling higher levels of software complexity and offering extensive functionality.
%
Remarkably, in our experimental evaluations, MetaGPT achieves a $100\%$ task completion rate, demonstrating the robustness and efficiency (time and token costs) of our design.

We summarize our contributions as follows:

\noindent \textbf{$\bullet$} 
We introduce MetaGPT, a meta-programming framework for multi-agent collaboration based on LLMs. It is highly convenient and flexible, with well-defined functions like role definition and message sharing, making it a useful platform for developing LLM-based multi-agent systems. 

\noindent \textbf{$\bullet$} 
Our innovative integration of human-like SOPs throughout MetaGPT's design significantly enhances its robustness, reducing unproductive collaboration among LLM-based agents.
Furthermore, we introduce a novel executive feedback mechanism that debugs and executes code during runtime, significantly elevating code generation quality (e.g., 5.4\% absolute improvement on MBPP).

\noindent \textbf{$\bullet$} 
We achieve state-of-the-art performance on HumanEval~\citep{chen2021evaluating} and MBPP~\citep{austin2021program}. 
Extensive results convincingly validate MetaGPT, suggesting that it is a promising meta-programming framework for developing LLM-based multi-agent systems.